\documentclass[12pt,a4paper]{article}
\usepackage{fontspec}
\IfFontExistsTF{Times New Roman}{
  \setmainfont{Times New Roman}
  \newfontfamily\cyrillicfont{Times New Roman}
}{
  \setmainfont{CMU Serif}
  \newfontfamily\cyrillicfont{CMU Serif}
}
\IfFontExistsTF{Courier New}{
  \setmonofont{Courier New}
  \newfontfamily\cyrillicfonttt{Courier New}
}{
  \setmonofont{CMU Typewriter Text}
  \newfontfamily\cyrillicfonttt{CMU Typewriter Text}
}

\usepackage{polyglossia}
\setdefaultlanguage{russian}
\setotherlanguage{english}

\usepackage[a4paper,left=20mm,right=16mm,top=20mm,bottom=20mm]{geometry}
\usepackage{setspace}
\onehalfspacing

\usepackage{indentfirst}
\setlength{\parindent}{1.25cm}
\usepackage[hidelinks]{hyperref}
\usepackage{xurl}
\emergencystretch=3em
\newcommand{\code}[1]{\texttt{\nolinkurl{#1}}}

\newcommand{\university}{Государственный Университет Молдовы}
\newcommand{\faculty}{Факультет Математики и Информатики}
\newcommand{\department}{Кафедра Информатики}
\newcommand{\worktype}{Проектная работа}
\newcommand{\projecttitle}{Сайт недвижимости UrbanNest Estate}
\newcommand{\studentname}{Splavski Maxim\\Metelska Daniela}
\newcommand{\studentgroup}{I2402}
\newcommand{\teacher}{Visnevschi V.}
\newcommand{\cityyear}{Кишинэу, 2026}

\newcommand{\makeprojecttitle}{
\begin{titlepage}
\thispagestyle{empty}
\begin{center}
\university\\
\faculty\\
\department

\vspace{4cm}
{\bfseries\fontsize{16pt}{19pt}\selectfont \worktype\par}
\vspace{0.8cm}
{\fontsize{14pt}{17pt}\selectfont Тема:\par}
\vspace{0.25cm}
{\bfseries\fontsize{15pt}{18pt}\selectfont \projecttitle\par}
\end{center}

\vspace{2.7cm}
\begin{flushright}
\begin{minipage}{0.45\textwidth}
Выполнили студенты: \textbf{\studentname}\\
Группа: \textbf{\studentgroup}\\[0.5cm]
Проверил преподаватель: \textbf{\teacher}
\end{minipage}
\end{flushright}

\vfill
\begin{center}
\cityyear
\end{center}
\end{titlepage}
\setcounter{page}{2}
}


\begin{document}
\makeprojecttitle
\tableofcontents
\newpage

\section{Теоретическая часть}

Веб-приложение UrbanNest Estate реализовано как традиционный PHP-проект, работающий через Apache и MySQL в среде XAMPP. Серверная часть отвечает за обработку форм, работу с базой данных, проверку пользователей и формирование HTML-страниц.

Архитектурно проект наследует принципы MVC. Модели отвечают за работу с таблицами базы данных, контроллеры обрабатывают запросы, а представления формируют интерфейс. Такой подход делает код понятнее и уменьшает смешивание логики с HTML.

Проектная работа выполнена двумя студентами: Splavski Maxim и Metelska Daniela. Оба участника указаны на титульном листе отчета как авторы работы.

\section{Формулировка задачи}

Необходимо разработать веб-приложение для просмотра и подбора современной недвижимости. Пользователь должен иметь возможность просматривать каталог, искать объекты, регистрироваться, входить в систему и отправлять заявку на просмотр. Администратор должен управлять объектами, типами недвижимости, заявками и пользователями.

\section{Цель и этапы работы}

Цель работы - создать завершенный MVP сервиса UrbanNest Estate на PHP и MySQL в рамках общей проектной работы двух студентов.

Основные этапы:
\begin{enumerate}
  \item проектирование структуры приложения;
  \item создание базы данных;
  \item реализация регистрации и входа;
  \item реализация каталога недвижимости;
  \item реализация формы заявки на просмотр;
  \item разработка административной панели;
  \item оформление интерфейса и документации.
\end{enumerate}

\section{Практическая часть}

Проект реализован без REST API и без внешних backend-фреймворков. Единой точкой входа является файл \code{public/index.php}. Он определяет маршрут и вызывает соответствующий контроллер.

Код разделен на директории \code{app/models}, \code{app/controllers}, \code{app/views}, \code{app/core} и \code{app/config}. Подключение к базе данных вынесено в отдельный класс \texttt{Database}, который создает PDO-соединение.

\section{Краткое описание особенностей реализации}

Приложение использует модульную MVC-подобную архитектуру. В директории \texttt{models} находятся классы для работы с пользователями, объектами недвижимости, типами недвижимости и заявками. Контроллеры принимают запросы, проверяют права доступа и передают данные в представления.

Маршрутизация реализована в файле \code{public/index.php}. Для Apache добавлен файл \code{public/.htaccess}, который направляет пользовательские URL в единую точку входа. Основные настройки проекта и базы данных находятся в \code{app/config/config.php}.

Фотографии объектов недвижимости хранятся локально. Для тестовых данных используются квадратные JPG-файлы размером 900 на 900 пикселей из \code{public/assets/img/properties}. Если администратор загружает новое изображение, система проверяет, что исходный файл имеет размер не меньше 900 на 900 пикселей, выполняет центральную квадратную обрезку, уменьшает результат до 900 на 900 пикселей и сохраняет готовый JPG-файл в \code{public/uploads/properties}. В базе данных сохраняется только относительный путь.

\section{Функциональные возможности}

В приложении реализованы:
\begin{itemize}
  \item главная страница с hero-блоком, быстрым поиском и избранными объектами;
  \item каталог недвижимости с поиском, фильтрами и сортировкой;
  \item детальная страница объекта;
  \item регистрация, вход и выход;
  \item личный кабинет пользователя;
  \item форма заявки на просмотр;
  \item административная панель;
  \item управление объектами, типами недвижимости, заявками и пользователями.
\end{itemize}

\section{Основные страницы приложения}

Главная страница содержит крупный заголовок, краткое описание сервиса, быстрый поиск и блок объектов недели. Каталог недвижимости выводит данные из базы и поддерживает поиск по названию, городу или району, фильтр по типу объекта, типу сделки, комнатам, статусу, цене и площади.

Страница объекта показывает изображение, город, район, адрес, цену, площадь, количество комнат, этаж, тип объекта, тип сделки, описание и статус. Форма заявки доступна авторизованным пользователям и содержит объект, имя, телефон, email, удобное время связи и комментарий.

Административная панель содержит dashboard со статистикой, управление объектами недвижимости, типами недвижимости, заявками и пользователями. Администратор может создавать новых администраторов и изменять роли пользователей.

\section{Сценарии взаимодействия}

Гость открывает сайт, просматривает каталог, фильтрует объекты по типу сделки и бюджету, затем изучает страницу выбранного объекта. Для отправки заявки он регистрируется или входит в аккаунт.

Пользователь выбирает объект недвижимости, указывает контактные данные, удобное время связи и комментарий. После отправки заявка появляется в личном кабинете.

Администратор входит в систему, открывает dashboard, просматривает статистику, добавляет объекты, редактирует типы недвижимости и меняет статусы заявок.

\section{Механизм регистрации, входа и ролей}

Регистрация выполняется через форму с именем, email, телефоном и паролем. На сервере проверяются длина имени, корректность email, телефон, длина пароля и уникальность email. Пароль сохраняется в таблице \texttt{users} только после обработки функцией \code{password_hash()}.

При входе пользователь вводит email и пароль. Контроллер находит пользователя по email и проверяет пароль функцией \code{password_verify()}. После успешного входа основные данные пользователя сохраняются в \code{$_SESSION}. Для доступа к личному кабинету используется проверка авторизации, а для административных страниц дополнительно проверяется роль \texttt{admin}.

\section{Структура базы данных}

База данных содержит четыре основные таблицы:
\begin{itemize}
  \item \texttt{users} - пользователи и роли;
  \item \texttt{property\_types} - типы недвижимости;
  \item \texttt{properties} - характеристики объектов недвижимости;
  \item \code{property_requests} - заявки пользователей на просмотр.
\end{itemize}

Таблицы связаны внешними ключами. Объект относится к типу недвижимости, а заявка связана с пользователем и объектом недвижимости.

\section{Формы и валидация}

В проекте реализованы несколько HTML-форм. Форма регистрации создает нового пользователя. Форма поиска в каталоге позволяет фильтровать объекты по нескольким критериям. Форма заявки содержит объект недвижимости, имя, телефон, email, удобное время связи и комментарий.

Административная форма объекта содержит более пяти полей: название, тип, сделку, город, район, адрес, комнаты, площадь, этаж, общую этажность, цену, статус, изображение и описание. На клиенте используются атрибуты \texttt{required}, \texttt{minlength}, \texttt{type=email}, \texttt{type=number} и другие. На сервере контроллеры выполняют дополнительную проверку данных перед сохранением.

\section{Безопасность}

Пароли сохраняются в базе только в хешированном виде. Для проверки пароля используется функция \code{password_verify()}. Доступ к личному кабинету защищен PHP-сессиями, а административный раздел дополнительно проверяет роль пользователя.

Все SQL-запросы выполняются через подготовленные запросы PDO. Данные при выводе экранируются функцией \code{htmlspecialchars()}, чтобы снизить риск XSS.

Загрузка изображений ограничена исходными форматами JPG, PNG и WEBP. Минимальный размер исходной фотографии - 900 на 900 пикселей. После загрузки изображение автоматически приводится к квадратному формату 900 на 900 пикселей и сохраняется в отдельную директорию \code{public/uploads/properties}, а имя файла генерируется автоматически.

\section{Выводы}

В ходе совместной работы было создано завершенное PHP + MySQL приложение UrbanNest Estate. Проект демонстрирует работу с базой данных, формами, ролями, сессиями, CRUD-операциями, фильтрацией и адаптивным интерфейсом.

Ссылка на репозиторий добавляется после загрузки проекта на GitHub.

\section{Ответы на контрольные вопросы}

\begin{enumerate}
  \item Серверная часть обрабатывает запросы, работает с базой данных и формирует страницы.
  \item База данных нужна для хранения пользователей, объектов недвижимости, типов и заявок.
  \item Аутентификация проверяет личность пользователя.
  \item Авторизация определяет доступные пользователю действия.
  \item Пароли нельзя хранить открыто из-за риска утечки.
  \item PHP-сессии сохраняют состояние входа пользователя.
  \item SQL-инъекция - внедрение вредоносного SQL через ввод пользователя.
  \item Prepared statements отделяют SQL-код от данных.
  \item Валидация проверяет корректность введенных данных.
  \item Админ-панель нужна для управления данными проекта.
  \item Роль пользователя определяет уровень доступа.
  \item Адаптивная верстка делает сайт удобным на разных устройствах.
  \item MVC разделяет данные, обработку и представление.
  \item Модульная структура легче поддерживается и проверяется.
\end{enumerate}

\section{Список использованных источников}

\begin{itemize}
  \item Документация PHP: \texttt{https://www.php.net/}
  \item Документация MySQL: \texttt{https://dev.mysql.com/doc/}
  \item Документация XAMPP: \texttt{https://www.apachefriends.org/}
  \item Wikimedia Commons: фотографии недвижимости для локальных тестовых данных.
\end{itemize}

\end{document}
